\section{Factores $R$ y $R_0$}

Para determinar los factores $R_0$ y $R$ correspondientes al sistema estructural 
del edificio, se recurre a la Tabla 5.1 de la norma NCh433~\cite{NCh433:1996}, la 
cual especifica dichos valores en función del material y del sistema resistente a 
las cargas sísmicas. Dado que el sistema estructural corresponde a muros de corte 
de hormigón armado, los factores a utilizar son:

\begin{equation}
    R = 7 \qquad R_0 = 11
\end{equation}

% ---------------------------------------------------------------
\section{Coeficientes Sísmicos Máximo y Mínimo}

Previo al análisis sísmico, la norma NCh433~\cite{NCh433:1996} establece rangos 
admisibles para el coeficiente sísmico $C$. A continuación se determinan el valor 
mínimo y máximo permitidos.

\subsection{Coeficiente Sísmico Mínimo}

Según el artículo 6.2.3.1 de la norma NCh433~\cite{NCh433:1996}, el coeficiente 
sísmico no puede ser inferior a:

\begin{equation}
    C_{\min} = \frac{A_0 \cdot S \cdot I}{6} = \frac{0{,}40 \cdot 0{,}9 \cdot 1{,}0}{6} 
    = 0{,}060
\end{equation}

\subsection{Coeficiente Sísmico Máximo}

Según el artículo 6.2.3.2 de la norma NCh433~\cite{NCh433:1996}, el coeficiente 
sísmico máximo se determina como:

\begin{equation}
    C_{\max} = 0{,}35 \cdot S \cdot A_0 \cdot I 
    = 0{,}35 \cdot 0{,}9 \cdot 0{,}40 \cdot 1{,}0 = 0{,}126
\end{equation}

Por lo tanto, el coeficiente sísmico de diseño debe satisfacer:
\begin{equation}
    0{,}060 \leq C \leq 0{,}126
\end{equation}

El valor de $C = 0{,}10$ se encuentra dentro del rango admisible, siendo consistente 
con lo indicado en la Tabla~\ref{tab:parametros_sismicos}, y será el utilizado en 
el análisis sísmico.

% ---------------------------------------------------------------
\section{Cálculo del Factor de Corrección $R^*$ para Sismo en X e Y}

Con los valores determinados en las secciones anteriores, se procede a calcular el 
factor de reducción de la fuerza sísmica $R^*$, mediante la expresión (6-10) de la 
norma NCh433~\cite{NCh433:1996}:

\begin{equation}
    R^* = \frac{T^*}{0{,}1 \cdot T_0 + \dfrac{T^*}{R_0}}
\end{equation}

donde $T^*$ corresponde al período del modo con mayor masa traslacional equivalente, 
y $T_0$ es el período característico del suelo definido en la 
Tabla~\ref{tab:parametros_sismicos}.

Para determinar $T^*$, se ejecutó el modelo en ETABS, obteniéndose los períodos de 
vibración indicados en la Tabla~\ref{tab:periodos}.

\begin{table}[H]
    \centering
    \caption{Períodos de los modos de vibración obtenidos en ETABS.}
    \label{tab:periodos}
    \begin{tabular}{ccc}
        \toprule
        \textbf{Modo} & \textbf{Eje traslacional} & \textbf{Período (s)} \\
        \midrule
        1 & Torsional & $1{,}251$ \\
        2 & Eje X     & $1{,}140$ \\
        3 & Eje Y     & $0{,}853$ \\
        \bottomrule
    \end{tabular}
\end{table}

Para el análisis se consideran los modos 2 y 3, correspondientes a los ejes X e Y 
respectivamente. El $R^*$ se calcula para cada dirección:

\begin{itemize}
    \item \textbf{Modo 2 (Eje X):}
    \begin{equation}
        R^*_X = \frac{1{,}14}{0{,}1 \cdot 0{,}15 + \dfrac{1{,}14}{11}} = 10{,}609
    \end{equation}

    \item \textbf{Modo 3 (Eje Y):}
    \begin{equation}
        R^*_Y = \frac{0{,}853}{0{,}1 \cdot 0{,}15 + \dfrac{0{,}853}{11}} = 10{,}217
    \end{equation}
\end{itemize}

Dado que en ETABS el espectro se ingresa mediante el inverso de $R^*$, los 
factores a utilizar son:

\begin{equation}
    \frac{1}{R^*_X} = \frac{1}{10{,}609} = 0{,}0943 \qquad
    \frac{1}{R^*_Y} = \frac{1}{10{,}217} = 0{,}0979
\end{equation}

Con estos factores ingresados en los espectros respectivos, se obtienen los cortes 
basales de la Tabla~\ref{tab:corte_basal_inicial}.

\begin{table}[H]
    \centering
    \caption{Corte basal inicial por dirección.}
    \label{tab:corte_basal_inicial}
    \begin{tabular}{cc}
        \toprule
        \textbf{Dirección} & \textbf{Corte basal (tonf)} \\
        \midrule
        X & $160{,}34$ \\
        Y & $201{,}92$ \\
        \bottomrule
    \end{tabular}
\end{table}

Comparando con los cortes mínimos exigidos por NCh433~\cite{NCh433:1996}, se verifica 
que ambos valores se encuentran por debajo del mínimo requerido de $766{,}05\ \text{tonf}$. 
Por lo tanto, se ajusta el $R^*$ para cumplir con la norma, utilizando la siguiente 
expresión:

\begin{equation}
    R^*_{\text{ajustado}} = \frac{Q_0}{Q_{\text{norma}}} \cdot R^*_{\text{calculado}}
\end{equation}

donde $Q_0$ es el corte basal obtenido del modelo, $Q_{\text{norma}}$ es el corte 
basal mínimo exigido por la norma y $R^*_{\text{calculado}}$ es el factor determinado 
previamente. Aplicando para cada dirección:

\begin{itemize}
    \item \textbf{Eje X:}
    \begin{equation}
        R^*_X = \frac{160{,}34}{766{,}05} \cdot 10{,}609 = 2{,}22
    \end{equation}

    \item \textbf{Eje Y:}
    \begin{equation}
        R^*_Y = \frac{201{,}92}{766{,}05} \cdot 10{,}217 = 2{,}69
    \end{equation}
\end{itemize}

Ingresando los nuevos factores en ETABS, se obtienen los cortes basales definitivos 
de la Tabla~\ref{tab:corte_basal_final}, los cuales cumplen con los requisitos 
establecidos en NCh433~\cite{NCh433:1996}.

\begin{table}[H]
    \centering
    \caption{Corte basal definitivo por dirección con $R^*$ ajustado.}
    \label{tab:corte_basal_final}
    \begin{tabular}{cc}
        \toprule
        \textbf{Dirección} & \textbf{Corte basal (tonf)} \\
        \midrule
        X & $766{,}06$ \\
        Y & $765{,}99$ \\
        \bottomrule
    \end{tabular}
\end{table}